\section{Results}\label{sec:results}

\subsection{RQ1: Does concurrency hurt standard RL, and does factored credit help?}
Table~\ref{tab:rq1} reports the insurer with 2, 3 and 6 independent components. Joint
PPO degrades as concurrency grows: its final return falls from 0.965 with one region to
0.748 with two, and with two regions it stays below 0.9 for all 40 epochs on 16 of 20
seeds. \nfgae{} holds its level (1.101 and 1.070): it is ahead of PPO at the end on 17
of 20 seeds with one region and on 20 of 20 with two, and over the whole run on 19 and
20 of 20 (paired $t$-test $p<0.002$ in every case). The gain grows with the number of
independent components (final $+0.034$, $+0.136$, $+0.321$ for $K=2,3,6$; whole-run
mean $+0.069$, $+0.218$, $+0.406$); with shared clerks the certificate merges claims
and complaints and the gain shrinks accordingly. Each process alone is learned by PPO
within 3--9 epochs (median over seeds of the first evaluation at or above 0.9); inside
the insurer PPO needs 18 with one region, and 3 of 20 seeds never get there, whereas
\nfgae{} needs 6. This locates the problem in the other processes' noise
(Fig.~\ref{fig:curves}). Fig.~\ref{fig:gaink} shows the gain against the number of
independent processes. On multi-site, where every method reaches the heuristic,
\nfgae{} gets there in half the epochs (4 vs 8, 20 of 20 seeds).

\begin{figure}[t]
\centering
\includegraphics[width=0.99\textwidth]{figures/fig_curves}
\caption{Insurer: greedy normalised return during training (mean and 95\% confidence
interval over 20 seeds; per-process: 10 seeds; 1 = heuristic). Per-process: separately
trained PPO agents, one per process, combined.}\label{fig:curves}
\end{figure}

\begin{figure}[t]
\centering
\includegraphics[width=0.55\textwidth]{figures/fig_gain_vs_k}
\caption{Gain of \nfgae{} over PPO against the number of independent processes $K$
(mean and 95\% confidence interval of the paired differences; 20 seeds, 10 for $K=2$):
$K=2$ shared clerks, $K=3$ one region, $K=6$ two regions. With $K=1$, \nfgae{} is PPO by
construction.}\label{fig:gaink}
\end{figure}

\begin{table}[t]
\centering
\caption{Insurer: normalised return (1 = heuristic), mean over 20 seeds (shared clerks
and per-process: 10). Per-process: one PPO agent per process, trained on its own
process's reward. With shared clerks,
claims and complaints form one coupled process, so there are no separate
per-process learners.}\label{tab:rq1}
\setlength{\tabcolsep}{4.6pt}
\begin{tabular}{lc ccc ccc}
\toprule
 & & \multicolumn{3}{c}{Final return} & \multicolumn{3}{c}{Mean over training} \\
\cmidrule(lr){3-5}\cmidrule(lr){6-8}
Variant & $K$ & PPO & Per-process & \nfgae & PPO & Per-process & \nfgae \\
\midrule
Shared clerks & 2 & 0.989 & --    & \textbf{1.023} & 0.855 & --    & \textbf{0.925} \\
1 region      & 3 & 0.965 & 1.058 & \textbf{1.101} & 0.811 & 0.960 & \textbf{1.029} \\
2 regions     & 6 & 0.748 & 1.048 & \textbf{1.070} & 0.611 & 0.950 & \textbf{1.017} \\
\bottomrule
\end{tabular}
\end{table}

\subsection{RQ2: When does waiting pay, and what does it take to learn it?}
In the slow-mismatch insurer, waiting for the matched underwriter is worth about 14\% of
underwriting's throughput. Table~\ref{tab:rq2} compares five arms with one and with two
regions. Only the combination of factored credit and component-scoped postponement
learns to wait reliably. With one region it reaches 1.054, at or above the waiting
heuristic on 18 of 20 seeds, and improves on \nfgae{} without postponement on every seed
($+0.127$). With two regions it reaches 1.005 (12 of 20 seeds at or above the waiting
heuristic) and improves on \nfgae{} without postponement by $+0.109$ (18 of 20 seeds,
Wilcoxon $p<0.001$): the value of waiting is about the same at both sizes. Giving PPO
the same option does not help: with one region its final return is not significantly
better than without ($+0.06$, $p=0.17$) and it is no better over the run ($-0.04$); with
two regions it is level with PPO without postponement ($+0.02$, 12 of 20 seeds,
$p=0.49$); the extra decisions cannot be credited through the other processes' noise.
Global postponement helps PPO a little with one region ($+0.09$ at the end, $p=0.011$,
nothing over the run) but not with two ($+0.02$, $p=0.40$): the more processes there
are, the more each global wait freezes. With two regions the gap between \nfgae{} with
component-scoped postponement and PPO's best arm grows from $+0.12$ to $+0.30$ (20 of 20
seeds, $p<10^{-5}$ by both tests).

The behaviour of the trained policies confirms the mechanism. \nfgae{} with
component-scoped postponement waits almost only in underwriting (5.8 times per episode
with one region, 10.7 with two; the waiting heuristic 6.5 and 15.3), almost never in
claims or complaints, where waiting hurts (0.07 times per episode with two regions), and
never while a useful assignment is available. Mismatched underwriting assignments fall
from 1.23 to 0.36 per episode with one region and from 2.71 to 1.17 with two. PPO's
waits with either kind of postponement include waits in claims and complaints while
useful work is available (1.2 and 1.9 per episode with two regions), and its mismatched
assignments do not fall. Fig.~\ref{fig:postpone} shows the final returns per seed and
the learning curves.

\begin{figure}[!tbp]
\centering
\includegraphics[width=\textwidth]{figures/fig_postpone}
\caption{Slow-mismatch insurer with one region (top) and two regions (bottom), 20 seeds.
Left: final normalised return per seed (dots) and mean (bar); 1 = waiting heuristic.
Right: greedy normalised return during training (mean and 95\% confidence
interval).}\label{fig:postpone}
\end{figure}

\begin{table}[t]
\centering
\caption{Slow-mismatch insurer, 20 seeds per arm. Return: normalised so that
1 = waiting heuristic (the work-conserving heuristic scores 0.873 with one region and
0.865 with two); ``$\geq 1$'': seeds that end at or above the waiting heuristic.
Behaviour: best checkpoint of each seed, averaged over 20 evaluation episodes, summed
over regions; the waiting heuristic waits 6.5 (one region) and 15.3 (two regions) times
per episode in underwriting.}\label{tab:rq2}
\setlength{\tabcolsep}{5pt}
\begin{tabular}{ll cc cc}
\toprule
 & & \multicolumn{2}{c}{Return} & \multicolumn{2}{c}{Behaviour per episode} \\
\cmidrule(lr){3-4}\cmidrule(lr){5-6}
Credit & Postponement & Final & $\geq 1$ & \shortstack{Waits in\\underwriting} & \shortstack{Mismatched\\assignments} \\
\midrule
\multicolumn{6}{l}{\emph{One region} ($K=3$)} \\
PPO    & none      & 0.847 & 0/20 & 0   & 1.52 \\
\nfgae & none      & 0.928 & 1/20 & 0   & 1.23 \\
PPO    & global    & 0.935 & 7/20 & 4.9 & 0.79 \\
PPO    & component & 0.910 & 7/20 & 4.4 & 1.12 \\
\nfgae & component & \textbf{1.054} & \textbf{18/20} & 5.8 & \textbf{0.36} \\
\midrule
\multicolumn{6}{l}{\emph{Two regions} ($K=6$)} \\
PPO    & none      & 0.685 & 0/20 & 0    & 3.91 \\
\nfgae & none      & 0.896 & 0/20 & 0    & 2.71 \\
PPO    & global    & 0.703 & 0/20 & 4.4  & 3.84 \\
PPO    & component & 0.707 & 0/20 & 3.7  & 4.19 \\
\nfgae & component & \textbf{1.005} & \textbf{12/20} & 10.7 & \textbf{1.17} \\
\bottomrule
\end{tabular}
\end{table}

\subsection{RQ3: Is it budget, tuning, or per-process learning?}
On the 2-region insurer, doubling PPO's data per epoch raises it from 0.719 to 0.856, and
tuning its learning rate to at most 0.782; \nfgae{} with half the data of the former
reaches 1.087, and at PPO's best learning rate 1.093 (seeds 0--4, ahead on 5 of 5 in
every comparison). Separate per-process learners capture most of the credit gain
(Table~\ref{tab:rq1}); \nfgae{} adds a shared network that learns faster early
(mean over epochs 3--12: 0.92 vs 0.77 at 2 regions) and a single policy, and, unlike
per-process learners, it can use component-scoped postponement under one policy.
\todo{Table: PPO variants}

\subsection{RQ4: Cost}
Encoding the whole organisation at every decision made \nfgae{} 1.7--2.2 times as
expensive as per-process learners, growing with the number of processes. Component turns
cancel this: on 2-epoch runs they speed \nfgae{} up 1.67 times with one region and 2.39
times with two, with unchanged learning (whole-run mean 1.023 vs 1.017, $p=0.73$).
\todo{Table: CPU time}

\subsection{The policy network}
On multi-site routing, PPO with a per-type HGT reaches 0.403 (6 of 20 seeds collapsed);
with the relational network of Sect.~\ref{sec:network}, the same PPO reaches 1.039. A
rerun of the HGT runs on the current code reproduces them bit for bit, so the difference
is the network. All comparisons above use the same network.
